\section{Discussion}

Our participants judged AI less by raw capability than by how well it fit the embodied, high-stakes conditions of their work. They found the most value in systems that eased cognitive load through memory, attention, and documentation support, but preferred these to remain passive in the background unless errors became critical. More broadly, this suggests a shift in framing; rather than designing AI as a collaborator that directly participates in scientific reasoning, \emph{AI may be more effective when treated as infrastructure that supports the conditions under which scientific reasoning happens.}


\subsection{Limitations of AI for Physical and Embodied Scientific Work}

Despite rapid advances in recent tools, AI still remains peripheral among our participants in the physically-oriented work of lab and field science. Corroborating the findings of prior studies~\cite{wagman2025generative, morris2023scientists}, it is used mostly `at the desk' for planning, literature review, data analysis, or writing. Despite the status quo, participants were able to use our speculative design activity to envision deeper, context-aware AI integration into their future workflows. This gap highlights structural mismatches, not necessarily technical limitations. As an analogy, current AI assistants resemble `generalist' programmers or technicians who can write boilerplate code but lack the bespoke domain expertise to identify which approaches are the most insightful within a scientific field. Current AI tools are also ill-suited for improvisational and materially grounded practices that characterize much of day-to-day scientific work `beyond the desk.'

While extensive prior work has posited that knowledge-based professions are more readily augmented by, or even displaced by, AI~\cite{tomlinson2025working, brynjolfsson2025canaries, maslej2025artificial, sherman2025amazon, de2023generative}, our study findings suggest that scientific domains involving physical manipulation, situated judgment, and tacit expertise currently resist such displacement of human labor. Tasks requiring bodily coordination and tool improvisation (e.g., meticulous pipetting techniques, DIY repairs in the field) exemplify the kinds of uncertain dynamic settings where scientists viewed human adaptability as irreplaceable. Iterative experimental work often requires scientists to learn new skills or obtain new materials and equipment. Automation needs therefore vary widely depending on individual situations and cannot be easily routinized with AI.

Our findings complement recent analyses on AI's impact on various occupations~\cite{tomlinson2025working}, which shows that roles involving manual labor, operating machinery, or other physically-oriented tasks currently have the least potential for AI collaboration. While Tomlinson et al.~\cite{tomlinson2025working} show that AI aligns best with knowledge work and communication-heavy tasks, our findings reveal how even within knowledge-intensive scientific roles, AI adoption is uneven, concentrated to desk-based activities, and largely absent from lab and fieldwork practices. Although we focused on \rev{scientific practitioners} working in physical domains, we view their work as a proxy for other occupations with similar core elements such as, say, hospital surgeons or emergency first responders.


\subsection{Future Opportunities for AI in Scientific Practice}

Speculative design activities at the end of each interview revealed directions for how AI could be integrated in ways that augment scientists' capabilities without displacing their judgment (\fig{fig:newspecdesign}). From our findings, we derive four higher-level design opportunities for AI, which address recurrent challenges that our participants faced. \rev{These opportunities build directly on the five speculative archetypes described in Section~\ref{sec:spec-design}.}
\\

\textbf{(1) Design to scaffold memory, attention, documentation, and knowledge transfer in physical tasks.} \rev{Archetype 1 (AI as a collective knowledge keeper) and Archetype 2 (AI as a distributed task monitor) show how participants struggle} with fragmented documentation and lapses in attention during complex embodied work, as well as difficulty capturing and understanding tacit knowledge. These \rev{archetypes} point toward designs of real-time, context-aware recording and summarization tools to reduce cognitive load `in the moment,' preserve expertise, and support knowledge transfer. \rev{Unlike desk-based tools that require users to stop and externalize knowledge, support should operate in the background during embodied activity. This distinction matters because physical workflows make it difficult to shift attention or focus away from the task at hand; knowledge must be captured and organized in situ rather than through retrospective, at-desk effort.}
\\

\textbf{(2) Tune AI proactivity to relative stakes.} \rev{Archetype 2 and Archetype 3 (AI as a monitor for scientists' cognitive/physical health) illustrate how participants determine interruptions by AI to be acceptable when errors or safety risks carry meaningful consequences, such as open incubator doors or unsafe physical or cognitive conditions.} \rev{Scientific practitioners} otherwise preferred systems that stayed passive, intervening only when errors had high stakes. Thus, enable AI to adapt its intervention level to task context, with user-configurable thresholds (e.g., use only visual cues in settings that are sensitive to sound, such as with animal subjects).
\\

\textbf{(3) Expand interaction modalities beyond the desk.} Arche-type 2, Archetype 4 (AI as a mobile scout for fieldwork), and Arche-type 5 (AI as a collaborator for physical chores) illustrate that scientists often work with their hands occupied, across multiple rooms and workspaces, or in remote outdoor settings where desktop/tablet interfaces are impractical. Participants envisioned solutions ranging from voice interfaces, large ambient displays, and physical indicators, to flying drones and modular robotic assistants. AI systems need to be embedded in and responsive to the environmental and spatial realities of lab and field work.
\\

\textbf{(4) Embrace digital naturalism in designs.} \rev{Archetype 4 highlights the need for AI systems that can adapt to local environments and are able to operate under mobility and limited infrastructure. Many future} AI designs could benefit from elements of \emph{digital naturalism}, a design philosophy that advocates for modular, field-adaptable, and open-ended scientific tools~\cite{makery2019quitmeyer}. Quitmeyer cautions against stripping away the rich, local contexts of scientific practice and environments~\cite{quitmeyer2024growing}. These tools should be materially embedded and co-configured on site, becoming extensions of scientists’ senses and mental models, and helping them form deeper connections to their work instead of abstracting it away ``at the desk.''

Taken together, these opportunities suggest designing AI as infrastructure integrated into scientists' sites of use.


\subsection{AI as Infrastructure instead of Collaborator: Supporting the conditions for human scientific reasoning}

% AI co-scientist citation: https://research.google/blog/accelerating-scientific-breakthroughs-with-an-ai-co-scientist/

Current public conversation around AI in science has focused on using it to automate reasoning and discovery, as exemplified by highly-funded corporate projects such as Google DeepMind's AlphaFold, which won the 2024 Nobel Prize for Chemistry~\cite{jumper2021highly}, along with follow-up systems like \emph{AI Co-scientist}, ``a multi-agent AI system [...] to accelerate the clock speed of scientific and biomedical discoveries''~\cite{gottweis2025towards}. Despite the attention given to these large-scale automation efforts, our study participants described the day-to-day of lab and field science as surprisingly ``artisan'' (in the words of P4), requiring hands-on skill, tacit knowledge, and situated improvisation while doing physical chores in the lab.

We saw that many challenges our participants faced were less about enhancing reasoning and \emph{more about sustaining the complex, tightly-coupled physical conditions that make deep scientific reasoning possible}. Participants derived fulfillment from discovery; it was the mental and procedural burdens of minimizing mistakes within their physical environments that they struggled with. Participants consequently did not envision AI as a collaborator in reasoning or decision-making. They instead imagined intelligent, context-sensitive \textbf{\textit{infrastructure}}--- tools that could passively monitor workflows, intervene to prevent errors, and reduce time wasted on unproductive troubleshooting attempts. 

The high stakes of error were surfaced as a major constraint to AI adoption. Mistakes meant weeks of lost labor or, in extreme cases, catastrophic consequences. With these conditions in mind, participants were open to AI systems that prevented lapses in attention or memory, but were skeptical towards AI tools that produced confident outputs without accountability or transparency. While pattern recognition and workflow monitoring were welcome, interpretive AI systems were seen as brittle, lacking the embodied, tacit knowledge that scientific judgment relies on.

These findings suggest a reframing of AI integration in scientific contexts. While much of recent work has focused on building AI systems that can reason, hypothesize, or automate decision-making, our study participants had different unmet needs. Instead, they imagined intelligent tools that strengthen the underlying infrastructures that make scientific reasoning possible: monitoring experiments, stepping in rapidly to prevent human lapses, which includes bodily limitations in memory, attention, and focus, and preserving context across their local communities, spaces, and time. This reframing aligns with prior critiques in human-centered computing, such as concerns that current computational systems risk imposing solutions that prioritize abstraction and \emph{detachment from physical context}---``reinventing virtually every other site of practice in its own image''~\cite{agre1997computation}. Suchman and Agre argue that all research practices, "even the most
analytic, [are] fundamentally concrete and embodied"~\cite{suchman1987plans}, and our findings reflect this. The \rev{scientific practitioners} we interviewed did not see reasoning as something that could be separated from the physical and practical conditions of their work, and were skeptical of AI tools that treated it that way. They advocated for more ``defensive'' speculative AI designs---systems that mitigate risk and preserve stability instead of maximizing speed and efficiency.

Despite advances in technology, many \rev{scientific practitioners} still rely on paper-based tools (e.g., P12’s laminated paper protocols, P8’s field notebooks), underscoring a persistent misalignment between current automation and scientists’ real-world needs. Two decades ago, Yeh et al. observed similar preferences among field biologists and created ButterflyNet, a hybrid physical–digital system to bridge this gap~\cite{yeh2006butterflynet}. Today, with AI and technology far surpassing 2006 capabilities, this is an opportune moment to revisit hybrid system designs grounded in the realities of scientific practice. 

In addition, the growing flexibility and declining costs of modern AI systems make them feasible for ad-hoc appropriation. This is especially relevant in the scrappy, resource-limited environments our participants described. In such settings, bespoke digital tools are costly, niche, and time-consuming to build, risking obsolescence as research directions shift. Fortunately, today’s AI models are becoming less expensive and easier to adapt, offering potential for AI tools that can be rapidly customized by end-users. While laboratory and field science provided the frame for our work, these findings might also offer insights for other high-stakes, embodied, and situated domains, including medical practitioners, artisanal craft workers, and field technicians that work in challenging physical settings.


